\section{Introduction}

Forecasting under partial observation requires more than choosing the model
with the smallest error in one table. A hospitalization-rate series can be
compatible with simple statistical persistence, autoregressive dynamics,
hand-specified compartmental assumptions, or structured search candidates with
different observation semantics. The practical question is therefore not only
which candidate wins, but which recommendation remains credible when evidence
is split across model families, selection objectives, proposers, and
evaluation modes.

This issue is acute in influenza hospitalization forecasting. Compartmental
models offer interpretable latent states and useful constraints, but
hospitalization rates are not direct observations of canonical infectious
prevalence. A model can fail because its latent dynamics are misspecified, or
because the observed quantity is a direct, delayed, mixed, or proxy view of the
latent process. Age strata may also differ in delay, susceptibility recycling,
and the relationship between observed hospitalization rates and infectious
burden. A credible recommendation must therefore use both agreement and
disagreement across evidence sources, while preventing leakage from held-out
test evidence and avoiding claims the evidence cannot support.

The central methodological question is:
\begin{quote}
How can we use agreement and disagreement across model families, proposers,
and evidence modes to decide which forecasting recommendations are credible
under partial observation, while preventing leakage and overclaiming?
\end{quote}

We answer this question with a verifier-gated, evidence-aware, cross-family
model-selection framework. The framework represents proposals as structured
candidate records, checks them against hard verifier rules, attaches
train/validation or rolling evidence, applies interchangeable policies, and
audits the boundary between supported and rejected claims. The same framework
can evaluate deterministic proposers, constrained API-assisted proposers,
random proposal orderings, failure-guided refinements, and oracle references
without allowing any proposer to generate executable model code or see
held-out test metrics. FluSurv-NET hospitalization forecasting is the real-data
case study, not a leaderboard claim: it compares forecasting baselines,
hand-specified epidemic models, constrained observation-aware structure
discovery, same-grammar ablations, bounded candidate execution, and compact
multi-season robustness.

The contributions are:
\begin{enumerate}[leftmargin=1.5em]
\item a verifier-gated cross-family model-selection framework for partially
observed forecasting candidates;
\item a constrained agentic proposer protocol with task schemas, explicit
allowlists, no-leakage contexts, JSON candidate outputs, and verifier
feedback;
\item a formal evidence and selection protocol with Pareto-epsilon, weighted
scoring, hard-veto policies, and claim-boundary auditing;
\item synthetic structured recovery and candidate-budget execution/replay
evaluations that test observation-label recovery, delay-label recovery, and
budget efficiency under controlled settings;
\item bounded real-data execution and a FluSurv-NET case study showing
age-dependent and objective-dependent recommendations; and
\item a four-provider structured proposer benchmark over OpenAI/GPT,
Anthropic/Claude, Google/Gemini, and DeepSeek under the same verifier and
no-leakage protocol.
\end{enumerate}

The empirical conclusions are deliberately bounded. The paper does not claim
state-of-the-art influenza forecasting, a FluSight leaderboard result, fully
autonomous science, unrestricted model discovery, or real-world mechanism
recovery. Instead, it shows how verifier-gated proposal, evidence comparison,
and claim auditing can turn heterogeneous evidence into cautious
recommendations. In the FluSurv-NET case study, the strongest positive
observation-aware discovery signal appears in the pediatric \texttt{0-4 yr}
series; adult strata often favor simpler forecasting baselines or
hand-specified epidemic models; and multi-season robustness remains mixed.

\paragraph{Relation to prior work.}
Operational influenza forecasting efforts, including collaborative FluSight
evaluations and ensemble systems, focus on prospective public-health
prediction and operational readiness \cite{reich2019collaborative,
reich2019accuracy,lutz2019pathforward,ray2024flusion}. Our FluSurv-NET study
is instead a retrospective case study for verifier-gated cross-family model
selection under partial observation, not a FluSight leaderboard entry.
Mechanistic influenza forecasting has shown the value of compartmental
structure and data assimilation for seasonal prediction
\cite{shaman2012forecasting,shaman2013realtime}; here, the emphasis is on
checking observation semantics, comparing against non-mechanistic baselines,
and auditing which claims the evidence supports. Symbolic-regression systems
such as SINDy, AI Feynman, and PySR search for interpretable equations from
data \cite{brunton2016sindy,udrescu2020aifeynman,cranmer2023pysr}; our
structured search is deliberately narrower and is embedded in a
verifier-gated selection layer rather than framed as unconstrained equation
discovery. Recent LLM/evaluator scientific-discovery systems motivate
proposal-and-evaluation loops \cite{romera2024funsearch,shojaee2024llmsr,
lu2024aiscientist}, but this work does not claim fully autonomous science:
API-assisted proposal is optional, allowlisted, and verifier-checked before
use. The evaluation design follows standard out-of-sample and rolling-origin
forecasting principles \cite{tashman2000outofsample,hyndman2021fpp3}, while
the main policy uses a Pareto-style multi-objective selection view inspired by
multi-objective optimization \cite{deb2002nsga2}.
